\chapter{Introduction}\label{chap:intro}
\section{Overview of React}
React (\href{https://react.dev/}{react.dev}) is the most widely used framework for developing web-based graphical user interfaces (GUI) today \citep{Sta24Stack}.
React developers structure the UI as a tree of \emph{components,} where each component serves as a \emph{functional} specification that declaratively defines a portion of the UI and its behavior.
To declare user interactions in a component, developers use \emph{Hooks}, a collection of library functions that ``hooks into'' the runtime to manage state and handle side effects within components.

A React component is a function that serves as a declarative specification of a UI, which is read by the React runtime to render the interface.
The process of ``reading'' this specification is nothing special---the runtime directly invokes the component, which is a plain function.
When these components are called with appropriate arguments, they return a data structure representing how the component should appear on the screen.
Ideally, when a component is called with the same arguments, it should render the same interface \citep{Cac16React}.

This functional approach differs significantly from traditional object-oriented UI frameworks.
In object-oriented frameworks like Flutter (\href{https://flutter.dev/}{flutter.dev}), Qt (\href{https://www.qt.io/}{qt.io}), UIKit (\href{https://developer.apple.com/documentation/uikit/}{developer.apple.com}), or even legacy versions of React \citep{Met25Component,MadLhoTip20Semantics}, components are modeled with classes, where each view directly corresponds to an instance of a class.
Developers define different methods to be called at different lifecycle phases---such as initialization, state update, etc.---creating a clear separation.

React, however, introduces a different model: a single function component plays multiple roles---the same component is called for both creating a new view and updating an existing view.
To enable this, React switches between different implementations of the same Hook at runtime depending on the rendering phase.

\begin{wrapfigure}[4]{l}{0.387\linewidth}
  \vspace*{-1.2\intextsep}
  \begin{reactcode}
function Counter({ x }) {
  const [s, setS] = useState(() => x);
  const h = () => setS((s) => s+1);
  return <button onClick={h}> {s}
    </button>; }
  \end{reactcode}
\end{wrapfigure}
\noindent The \reacttt{Counter} component on the left demonstrates the \reacttt{useState} Hook for managing state.
This simple component increments its state~\reacttt{s} each time the button is clicked.
When first rendered, \reacttt{useState} returns an initial value from the initializer function~\reacttt{() => x}.
Here,~\reacttt{s} captures the current state, and \reacttt{setS}~accepts an updater to set a new state.
In subsequent renders, \reacttt{useState} automagically ``remembers'' the previous state and returns it instead.
\reacttt{button} is correctly rendered with new~\reacttt{s}, by efficiently diffing the returned \reacttt{button}---this is called \emph{reconciliation} \citep{Met25Reconciliation}.
This example illustrates how even a simple component like \reacttt{Counter} has different semantics across lifecycle phases.

Another peculiarity is that developers \emph{cannot} use JavaScript (JS) variables and direct assignments to manage a component's state.
In the above example, assigning a new value to~\reacttt{s} does not update the state of~\reacttt{Counter}---\reacttt{setS} must be used to update the state.

\begin{wrapfigure}[8]{r}{0.44\linewidth}
  \vspace*{-.8\intextsep}
  \begin{reactcode}
function Cond() {
  const [b, setB] = useState(() => false);
  if (b) {
    const [s, setS] = useState(() => 0);
    return s;
  } else {
    const h = () => setB(b => !b);
    return <button onClick={h}> Show
      </button>; } }
  \end{reactcode}
\end{wrapfigure}
Moreover, developers must adhere to specific rules imposed by React, collectively known as the \emph{Rules of React} \citep{Met25Rules}.
These rules represent invariants that components must maintain for React to properly manage the UI.
The \reacttt{Cond} component on the right violates one of these rules on line~4 by calling \reacttt{useState} conditionally inside the \reacttt{if}~branch.
React requires that all Hooks be called unconditionally at the top-level of a component, ensuring they execute in the same order on every render.
Initially, \reacttt{Cond} will show a button to toggle the Boolean state~\reacttt{b} from \reacttt{false} to \reacttt{true}.
When you click ``Show'' (setting~\reacttt{b} to \reacttt{true}), instead of displaying the initial state of~\reacttt{s} which is~|0|, the UI breaks down at runtime.
This stems from React's implementation detail where each Hook call is identified by its position in a linked list of internal bookkeeping objects.


\section{The Problem}
Hooks have subtle semantics, making it challenging to avoid unexpected behaviors.
The rendering process is opaque---it is unclear \emph{how} React registers updates or \emph{when} re-renders occur.
Therefore, React developers must follow specific rules or risk breaking their UI (detailed in~\cref{chap:pit}).
Worse yet, the host language is oblivious to React's rules, so it's up to programmers to adhere to them.
Without a clear understanding of the underlying mechanisms, developers must rely on informal resources that typically provide high-level guides rather than explanations of the render semantics, leaving even experienced developers vulnerable to misconceptions about how Hooks actually work.


\section{Our Solution}
We present \Rtrace, an operational semantics of React Hooks:
\begin{itemize}
  \item \textbf{\Rtrace{} is a semantics for clarifying the behavior of React Hooks.}
    \Rtrace{} is \emph{high-level enough} to abstract away implementation details and \emph{precise enough} to aid developers in reasoning about rendering behavior~(\cref{sec:example}).
    Our semantics helps developers clearly understand confusing UI bugs like infinite re-rendering~(\cref{chap:pit}).

  \item \textbf{\Rtrace{} is a foundation for building semantic-based tools.}
    \Rtrace{} provides a semantics for React Hooks~(\cref{chap:sem}), which is the first step in designing any semantic-based tool, such as an abstract interpreter \citep{CouCou77Abstract,RivYi20Introduction}.
    As an example for a semantic-based tool based on \Rtrace, we present an interactive tool that explains and visualizes the behavior of React programs~(\cref{chap:impl}).

  \item \textbf{\Rtrace{} captures the essence of React.}
    We prove that \Rtrace{} conforms to the key properties of React according to the documentation, as well as presenting an empirical conformance test suite~(\cref{chap:rules}).
\end{itemize}


\paragraph{Organization}
We introduce \reacttt{useState} and \reacttt{useEffect} and describe challenges in understanding them in~\cref{chap:over}, and illustrate common pitfalls from their interactions in~\cref{chap:pit}.
We formalize React Hook semantics in~\cref{chap:sem}, present a \Rtrace-definitional interpreter and visualizer in~\cref{chap:impl}, and demonstrate conformance with React in~\cref{chap:rules}.
We discuss extensions and implications in~\cref{chap:discussion}, compare with related work in~\cref{chap:related}, and conclude with future work in~\cref{chap:conclusion}.
